% !TEX root = ../main.tex
%%%%%%%%%%%%%%%%%%%%%%%%%%%%%%%%%%%%%%%%%%%%%%%%%%%%%%%%%%%%%%%%%%%%%%%
% Conclusion
%%%%%%%%%%%%%%%%%%%%%%%%%%%%%%%%%%%%%%%%%%%%%%%%%%%%%%%%%%%%%%%%%%%%%%%

\section{Conclusion}
\label{sec:conclusion}

We have built a stock-flow-consistent agent-based economy in which investment
is endogenous, capital accumulates, labour is hired endogenously, and the wage
responds to unemployment --- and in which, consequently, the Kaleckian question
of whether retention expands or depresses output is an outcome rather than an
assumption.

Conditionally, that outcome has a clean structure --- though not the one a
single label would suggest. The response $Y(\rho)$ is U-shaped, and the
statement that survives every choice of estimator is about the turn rather
than the sign: \emph{retention depresses output while $\rho$ lies below a
$\sigma$-dependent threshold $\rhostar$ and expands it beyond, $\rhostar$
rises with $\sigma$, and the empirically anchored retention rate lies below
$\rhostar$ for every $\sigma \geq 0.5$}. Summarised through one slope, that
structure appears as a sign frontier at $\sigstar \approx 0.65$: a frontier
that \emph{rises} when wages are made flexible, because the paradox of costs
dominates substitution; that is invariant to expectation dynamics; and that a
balanced-budget unemployment benefit \emph{eliminates} entirely, making
retention expansionary at every elasticity by closing the leak through which
the unemployed drain the circular flow. We use ``wage-led'' and ``profit-led''
in that reduced sense throughout --- as local descriptions of one side of a
turn, declared as such --- and not as a verdict on the economy.

Marginally, that structure largely dissolves. Across the defensible parameter
space, wage-led outcomes are the exception, $\sigma$ explains about 2\% of the
variance, and survival is governed almost entirely by the depreciation rate. At
the depreciation rate the data imply, the model does not exist --- which, read
together with its inability to match $I/Y$ and $K/Y$ jointly, identifies a
single structural culprit: a capital block without trend growth.

Between those two readings sits a methodological result we did not expect to
report. Part of the apparent gap was the measuring instrument: a two-point
chord across a U-shaped response is not the whole-support slope that defined
the frontier, and replacing it moves the headline probability from $0.095$ to
$0.026$ while leaving the variance decomposition intact. And part of the gap
was never there at all --- an inversion of direction we had written into our
own summary, contradicted by our own tables.

The finding we would carry forward is neither of these. It is the
wage$\,\to\,$unemployment$\,\to\,$capital-erosion channel, which emerged
unlooked-for from the wage curve, killed three consecutive stabilisation
hypotheses, and proved immune to every demand-side instrument applied to it.
Demand policy works where demand binds. Where the wage bill drives an
oscillation that erodes the capital stock, more demand is not a stabiliser ---
and in the sharpest case measured here, a saturated fiscal instrument financed
by taxing workers to pay workers simply amplifies the cycle that kills the
economy.
